\chapter{Why This Work Exists}

\section{Start with the body}

Raise your arm and touch something in front of you. That felt like a single
decision. It was not. Somewhere between deciding and moving, your nervous system
settled dozens of separate questions: how hard should the biceps pull, how hard
the triceps, the three heads of the deltoid, and so on, renegotiated continuously
for the whole half-second the movement takes.

Knowing those individual muscle forces is genuinely useful:

\begin{itemize}[leftmargin=1.4em, itemsep=3pt]
\item \textbf{Prosthetics and exoskeletons.} A powered device that assists a
weakened arm must know what the arm is trying to do, and how hard, in order to
supply the difference.
\item \textbf{Rehabilitation.} After a stroke, recovery is often tracked by how
movement is \emph{produced}, not only whether the hand arrives. Two patients can
touch the same target with completely different muscle strategies, and the
difference is clinically meaningful.
\item \textbf{Surgery and injury.} Whether a tendon transfer will work, or
whether a joint is being overloaded, depends on force in individual muscles ---
not on the visible motion.
\end{itemize}

The difficulty is that muscle force is almost never measured. Putting a sensor on
a tendon inside a living person is not routinely possible. So the forces are
\emph{computed} from the movement, and that computation is the subject of this
work.

\section{The counting problem that makes it hard}

Here is the whole difficulty in one observation.

Our model of the arm has \textbf{four ways to move} --- the shoulder can swing
forward, out to the side, and rotate; the elbow can bend --- and \textbf{nine
muscles} to move them with. Nine controls, four things to control.

\begin{plain}
Imagine nine people holding ropes attached to a single heavy door, and you want
the door at one particular angle.

They could all pull gently together. Or three could pull hard while six rest. Or
--- wastefully --- four could pull one way while five pull the other, mostly
cancelling, everyone exhausted, the door in exactly the same place.

Stand across the room and watch only the door. All three look identical. The door
angle does not tell you who was pulling.
\end{plain}

That is the \textbf{muscle redundancy problem}, and it is not a limitation of our
instruments. It is a mathematical fact: with more muscles than degrees of
freedom, infinitely many force combinations produce exactly the same movement.
Watching the movement can never, even in principle, tell you the forces.

\begin{why}
This is why the thesis cannot be answered by simply measuring how well a
controller reaches its target. Reaching accuracy is the \emph{door angle}. Two
controllers with identical reaching accuracy can be using completely different
muscles. Everything in this project that looks like a detour --- the force
comparison, the co-contraction analysis, the synergy decomposition --- exists
because the obvious measurement cannot answer the question.
\end{why}

\section{The classical answer, and what it costs}

Since the movement does not determine the forces, an extra rule is needed to pick
one answer from the infinitely many. The rule used in biomechanics since
Crowninshield and Brand (1981) is \textbf{minimum effort}: among all combinations
that produce the required movement, choose the one where the muscles work least
hard.

Written out, the method --- called \textbf{static optimisation} --- has two steps.

\textbf{Step one, inverse dynamics.} From the recorded motion and the limb's mass
properties, compute the twisting force (torque) needed at each joint. This step
is exact. There is one answer.

\textbf{Step two, load sharing.} Distribute those joint torques among the muscles
according to the minimum-effort rule:
\[
\min_{f}\ \sum_{i=1}^{9}\left(\frac{f_i}{F_{\max,i}}\right)^{2}
\quad\text{subject to}\quad
R^{\top} f = \tau,\qquad 0 \le f_i \le F_{\max,i}.
\]

\begin{plain}
Reading that formula in words:

\emph{Choose the nine muscle forces $f$ that make the total effort as small as
possible} --- where each muscle's effort is its force as a fraction of its own
maximum strength, squared.

\emph{Subject to two conditions.} First, the forces must actually produce the
required joint torques $\tau$ (that is $R^{\top} f = \tau$, where $R$ holds each
muscle's leverage about each joint). Second, no muscle may pull harder than it
can, and none may push --- muscles only pull. That is $0 \le f_i \le F_{\max,i}$.

Why square the effort? Squaring punishes large values disproportionately, so
sharing work among several muscles beats making one muscle do everything. It also
makes two people pulling against each other very expensive, because both
contributions count and neither cancels in the cost.
\end{plain}

This method works and is standard. Its drawback is that it solves a fresh
optimisation problem \emph{at every instant of the movement}. That is fine for
analysing a recording afterwards. It is awkward inside a device that must respond
in milliseconds.

\section{Why anyone would try machine learning here}

A neural network has the opposite cost profile: expensive to train once, then
very cheap to use --- a fixed number of multiplications, identical for every
input, with no iteration and no convergence check.

\begin{plain}
Classical optimisation is like solving a puzzle from scratch every time you are
asked. A trained network is like having memorised the \emph{pattern} of answers:
it responds instantly, but only as well as its training taught it.

The obvious question is whether the memorised pattern is actually right --- and
that is exactly the question this thesis asks.
\end{plain}

\begin{trap}
The original motivation was speed, and speed turned out not to survive
measurement. The classical solve was timed at \SI{2165}{\micro\second} against
the network's \SI{295}{\micro\second}, a $7.3\times$ advantage --- but that timed
a \emph{different} calculation (the muscle-model inversion inside the simulator,
not the load-sharing problem), and both sides were unoptimised Python. Re-timed
properly, a direct solve of the actual load-sharing problem takes
\SI{24}{\micro\second}, about \textbf{ten times faster} than the network. The
speed argument was withdrawn. Chapter~\ref{ch:results} gives the numbers.
\end{trap}

\section{What this thesis actually asks}

\begin{center}
\begin{tabular}{@{}p{3pt}@{\hspace{10pt}}p{0.93\linewidth}@{}}
\cellcolor{navy} & \textbf{The question.} For a reaching movement of a
nine-muscle arm, does a controller trained by deep reinforcement learning produce
the \emph{same per-muscle forces} as classical static optimisation?
\end{tabular}
\end{center}

And, because the answer turned out to have two halves that behave differently:

\begin{enumerate}[leftmargin=1.5em, itemsep=3pt]
\item Does it recruit the \emph{same muscles in the same proportions}? ---
the \emph{pattern} of the solution.
\item Does it use the \emph{same amount of force}? --- the \emph{economy} of the
solution.
\end{enumerate}

The short answer, established in Chapter~\ref{ch:results}: the pattern agrees
well; the economy did not, until the objective was changed to ask for it, at
which point it largely did too.

\section{An important boundary on what can be claimed}

One limitation shapes everything and should be understood before any result is
read.

Both the learned controller and the classical reference are computed inside the
same physics simulator, against the same model of muscle. Where they agree, two
computations over the same approximation agree. \textbf{No real arm, no
electromyography, and no human subject enters this work at any point.}

Furthermore, static optimisation is itself a \emph{theory} of how the nervous
system shares load --- a well-established and useful one, but contested. A
controller that disagrees with it is not thereby shown to be unphysiological.

\begin{why}
This is stated first rather than buried in a limitations section because it sets
the meaning of every subsequent number. ``Agreement of $r = 0.94$ with the
classical solution'' means \emph{this controller resembles the standard
computational method}. It does not mean \emph{this controller resembles a human}.
Those are different claims and only the first is supported.
\end{why}
